\section{Conclusions}
\label{sec:conclusions}

This paper introduced CascadeBench, a black-box benchmark that measures how
acoustic channel degradations propagate through commercial speech
recognition and translation services, and reported a complete measurement
over a controlled bilingual corpus of 480 sentences, fourteen degradation
conditions from seven families, and an eight-level noise severity sweep
from 20 to $-10$~dB, covering five recognition
systems from three product classes, three cascade translation paths,
and two end-to-end speech translation paths, all accessed through their
public interfaces on free service quotas; the central design was validated
on 150 human-recorded utterances under measured impulse responses and a
six-level real-babble ladder, a measured-RIR bridge condition on the
synthetic corpus, an open-weights locally executed recogniser, and a
second-vendor translation engine used as both a family-disjoint reference
and an alternative cascade backend.

The measurements support five conclusions. First, the product classes of
commercial recognisers differ in where they fail under noise, not merely how
fast: LLM-based recognisers hold a sub-percent floor and cross ten percent
error 2.2 to 2.5~dB later in Mandarin and 4.4 to 4.8~dB later in English
than the specialised streaming system, then collapse more sharply, and below
0~dB they fail by insertion-heavy hallucination with character error rates
above 100 percent; the real-speech validation reproduces this
knee-and-collapse structure and the hallucination regime, with three
calibrations: real rooms at matched reverberation time prove far milder
than the synthetic low-DRR stress response (and a measured telephone-band
specialist loses 3.4 rather than 18.2 percent under them), the LLM
noise-tolerance advantage holds only in the near-knee, low-error regime --
disappearing and for one LLM system reversing once error passes the
ten-percent level -- and is larger against real babble than synthetic
babble, and the Mandarin-insertion, English-deletion failure structure
replicates on an open-weights model outside the vendor, while
language-specific tolerance levels do not. Second, robustness is
condition- and language-specific
rather than uniform: narrowbanding, coding, packet loss, clipping, and tempo
changes at deployed severities are individually negligible for all systems,
0~dB babble nearly destroys English recognition for every system while
leaving the best Mandarin recognition low, and no system, of any class,
transcribes reliably below 0~dB. Third, the quality of a speech translation
cascade under degradation is governed by its recognition frontend, and the
relationship is an empirical linear pass-through of 0.44 to 0.54
BLEU points per percentage point of recognition error, with intercepts
statistically indistinguishable from zero, stable under leave-one-condition-out
refits for two of the three cascades, and honest condition-level
uncertainty that is substantially wider than the sentence-level intervals,
across up to fifteen conditions and three frontends; on real speech the
same relation holds with slopes of 0.43 (Mandarin) and 0.72 (English) over
nine conditions, and it survives swapping the translation engine for a
second vendor's model, which moves absolute scores by tens of BLEU points
but leaves the relative-loss trajectory nearly unchanged. Fourth,
end-to-end speech translation does not inherit
robustness automatically: the most resilient translation path measured is a
dedicated streaming end-to-end service, while the omni-modal end-to-end
path is the slowest and the most fragile of the end-to-end paths, degrading
more than the LLM-frontend cascades at 0~dB in point estimate under every
reference, with a margin that shrinks to under one point against the
Qwen3-ASR cascade under the oracle reference and whose paired interval
crosses zero there; on real speech the streaming end-to-end path likewise
loses a larger fraction of its clean quality under 0~dB real babble than a
well-frontended cascade, while the absolute BLEU ordering of the two paths
remains reference-dependent. Fifth,
black-box evaluation carries specific methodological obligations that this
benchmark operationalises, including three-reference scoring with one
family-disjoint reference to control machine-reference affinity -- under
which the relative-loss ranking is stable except for a sub-point inversion
between the omni path and the Qwen3-ASR cascade under the same-engine
reference -- exact-transcript corpora to enable sub-percent
measurement, the reporting of response completeness alongside error rates,
and the anchoring of service-level findings in measured signal properties,
which here exposed a sharp dissociation between intelligibility metrics and
recognition error across degradation families.

The benchmark recipe, including the corpus, the seeded degradation pipeline,
the degraded audio tree, all model responses, and the analysis code, is
released so that the measurement can be re-executed and verified as
providers update their services. The released artefacts are versioned alongside the service
identifiers and access dates of every system, so that future re-runs can
attribute score movements either to provider-side model updates or to
changes in the benchmark itself, a distinction that periodic evaluation of
moving targets requires. Together with the zero-cost execution protocol
described in Section~\ref{sec:methods}, the release is intended to make
robustness re-measurement a routine maintenance activity for speech
deployments rather than a one-off research exercise. As
speech intelligence continues to migrate into provider-operated black-box
systems, controlled measurement of this kind, repeated over time and across
providers, offers the community a practical substitute for the training-set
access that classical robustness research assumed, and offers integrators
the quantitative basis that architectural intuition alone cannot supply.
